% 专题十二：存在与任意——否定、解集结构与存在性
\section{专题十二：存在与任意——否定、解集结构与存在性}

\begin{learnobjective}
\begin{itemize}
  \item 会把 ∃ 与 ∀ 互相否定：$\neg\exists\iff\forall\neg$，$\neg\forall\iff\exists\neg$；
  \item 会把「存在」「任意」翻译成解集是否非空、是否为定义域；
  \item 会用零点定理等存在性定理证明 ∃，不必找出那个具体的值；
  \item 会用反证法证明 ¬∃ 与 ¬∀，并知道哪一边通常更难；
  \item 会用长度、距离、个数等量把「结构」变成矛盾。
\end{itemize}
\end{learnobjective}

\begin{preview}
\subsection{导语：把量词变成集合，把集合变成矛盾}
「存在 $x$ 使 $P(x)$」与「对任意 $x$ 都有 $P(x)$」，是高中最常用的一对说法。
它们看着像反义词，其实可以互相否定、互相翻译：
把量词换成解集，把解集换成数轴上的结构，再找一个可量化的量（长度、距离、个数），
往往就能造出矛盾。
本专题用六个视角把三个“种子问题”各做一遍。

\subsection{种子问题}
先自己做，再看后面的六个视角。
\begin{enumerate}
  \item 证明：存在 $x\in(1,2)$，使 $x^3-x-1=0$。
  \item 证明：对任意 $x\in\mathbb{R}$，都有 $x^2+1>0$。
  \item 证明：不存在长度大于 $0$ 的区间 $I$，使 $\sin x=0$ 对所有 $x\in I$ 成立。
\end{enumerate}
\end{preview}

\begin{explore}

\subsection{视角一：否定——∃ 与 ∀ 的互化}

\key{要点}
\begin{itemize}
  \item $\neg(\exists x\,P(x))\iff\forall x\,\neg P(x)$；
        $\neg(\forall x\,P(x))\iff\exists x\,\neg P(x)$。
  \item 否定量词时，量词要换（∃↔∀），结论也要否定。
  \item 双重否定：$\neg\exists x\,\neg P(x)\iff\forall x\,P(x)$（「不存在不满足的」=「全都满足」）。
\end{itemize}

\begin{example}
把「方程 $x^2+1=0$ 无实数解」写成符号形式，并写出它的否定。
\end{example}
\begin{solution}
无解即 $\forall x\in\mathbb{R},\ x^2+1\neq0$，也可写成 $\neg\exists x\in\mathbb{R},\ x^2+1=0$。
它的否定是 $\exists x\in\mathbb{R},\ x^2+1=0$。
\key{（回头看：种子问题 2 就是这条的正面）}
\end{solution}

\begin{problem}
写出「存在 $x\in\mathbb{R}$，使 $f(x)>0$」的否定，并说明它的意思。
\end{problem}
\begin{solution}
否定为 $\forall x\in\mathbb{R},\ f(x)\leq0$，即「$f$ 处处不大于 $0$」。
\end{solution}

\subsection{视角二：解集的结构——把条件变成集合}

\key{要点}
\begin{itemize}
  \item $\exists x\,P(x)\iff S=\{x\mid P(x)\}\neq\varnothing$（解集非空）；
  \item $\forall x\,P(x)\iff S=D$（解集等于定义域）；
  \item $\neg\exists\iff S=\varnothing$；$\neg\forall\iff S\subsetneq D$（有反例）。
  \item 结构上看：解集是区间、离散点集、还是空集。
\end{itemize}

\begin{example}
设 $D=\mathbb{R}$，$S=\{x\mid x^2-2x-3<0\}$。写出「$\forall x\in S,\ x<4$」的集合含义，并判断。
\end{example}
\begin{solution}
$S=\{x\mid(x-3)(x+1)<0\}=(-1,3)$。
「$\forall x\in S,\ x<4$」即 $S\subseteq(-\infty,4)$，显然成立。
\key{（回头看：把 ∀ 翻译成集合的包含，正是集合语言的用处）}
\end{solution}

\begin{example}
设 $S$ 满足：$1\in S$，且若 $a\in S$ 则 $a+1\in S$。写出 $S$ 的结构。
\end{example}
\begin{solution}
由生成规则，$S=\{1,2,3,\dots\}$。它不是「随便一个集合」，而是被规则完全确定。
\key{（有的集合要靠规则读出结构，而不是靠列举）}
\end{solution}

\begin{problem}
把「存在 $x\in[0,2]$ 使 $x^2-3x+2<0$」写成集合语言，并判断真假。
\end{problem}
\begin{solution}
解集 $S=\{x\mid(x-1)(x-2)<0\}=(1,2)$，非空，故命题为真。
\end{solution}

\subsection{视角三：存在性定理——白送的 ∃}

\key{要点}
\begin{itemize}
  \item 零点定理：$f$ 在 $[a,b]$ 连续且 $f(a)f(b)<0$，则 $\exists c\in(a,b)$ 使 $f(c)=0$。
  \item 介值定理、最值定理、中值定理、抽屉原理，都能白送「存在」。
  \item 用定理证 ∃，不需要把 $c$ 求出来。
\end{itemize}

\begin{example}
证明：存在 $x\in(1,2)$，使 $x^3-x-1=0$。
\end{example}
\begin{solution}
令 $f(x)=x^3-x-1$。$f(1)=-1<0$，$f(2)=5>0$，由零点定理，
$\exists c\in(1,2)$ 使 $f(c)=0$。
\key{（回头看：种子问题 1。注意这里根本没算 $c$ 等于几）}
\end{solution}

\begin{example}
证明：$13$ 个人中，至少有 $2$ 个人的生日在同一个月。
\end{example}
\begin{solution}
把 $13$ 个人放进 $12$ 个月，由抽屉原理，必有一个月至少有 $2$ 个人。
\key{（存在性定理的威力：只靠计数，不靠具体排列）}
\end{solution}

\begin{problem}
证明：存在 $x\in(0,1)$，使 $x^3+x-1=0$。
\end{problem}
\begin{solution}
令 $f(x)=x^3+x-1$，$f(0)=-1<0$，$f(1)=1>0$，由零点定理得证。
\end{solution}

\subsection{视角四：反证法——证明 ¬∃ 与 ¬∀}

\key{要点}
  \begin{itemize}
  \item 证 ¬∃（不存在 / 无解）：先假设 ∃，再推出矛盾。
  \item 证 ¬∀（有反例）：直接找反例。
  \item 反证法的三步：否定结论 → 推出矛盾 → 肯定原结论。
  \item 矛盾可以来自已知条件、公理、定理，或假设本身。
\end{itemize}

\begin{example}
证明：对任意 $x\in\mathbb{R}$，都有 $x^2+1>0$。
\end{example}
\begin{solution}
反证：假设 $\exists x\in\mathbb{R}$ 使 $x^2+1\leq0$，则 $x^2\leq-1$，
与 $x^2\geq0$ 矛盾。故原命题成立。
\key{（回头看：种子问题 2）}
\end{solution}

\begin{example}
证明：$\sqrt2$ 不是有理数。
\end{example}
\begin{solution}
反证：设 $\sqrt2=\dfrac pq$（$p,q$ 互质），则 $p^2=2q^2$，故 $p$ 为偶数。
设 $p=2k$，代入得 $2k^2=q^2$，故 $q$ 也为偶数，与 $p,q$ 互质矛盾。
\end{solution}

\begin{problem}
用反证法证明：对任意 $x\in\mathbb{R}$，都有 $x^2+2x+2>0$。
\end{problem}
\begin{solution}
假设 $\exists x$ 使 $x^2+2x+2\leq0$，即 $(x+1)^2+1\leq0$，得 $(x+1)^2\leq-1$，矛盾。
\end{solution}

\subsection{视角五：长度、距离与个数——把结构变成矛盾}

\key{要点}
\begin{itemize}
  \item 长度大于 $0$ 的集合一定非空：长度是「有 ∃」的便宜办法。
  \item 覆盖：若有限个区间覆盖 $[0,1]$，则长度之和至少为 $1$。
  \item 离散点集不含区间：区间长度 $>0$，点集长度 $=0$。
  \item 距离：两个不同点的距离 $>0$，可用来造「多出一个点」的矛盾。
\end{itemize}

\begin{example}
证明：长度总和小于 $1$ 的有限个开区间，不能覆盖 $[0,1]$。
\end{example}
\begin{solution}
反证：设 $[0,1]\subseteq\bigcup_{k=1}^n I_k$。则
\[
  1=|[0,1]|\leq\left|\bigcup_{k=1}^n I_k\right|\leq\sum_{k=1}^n|I_k|<1,
\]
矛盾。
\key{（回头看：把「覆盖」翻译成长度，结构就变成了不等式）}
\end{solution}

\begin{example}
证明：不存在长度大于 $0$ 的区间 $I$，使 $\sin x=0$ 对所有 $x\in I$ 成立。
\end{example}
\begin{solution}
$\sin x=0$ 的解集是 $\{k\pi\mid k\in\mathbb{Z}\}$，由离散的点组成，长度为 $0$。
任何长度大于 $0$ 的区间都不可能含在它里面，故不存在。
\key{（回头看：种子问题 3）}
\end{solution}

\begin{problem}
证明：不存在长度大于 $0$ 的区间 $I$，使 $\cos x=2$ 对所有 $x\in I$ 成立。
\end{problem}
\begin{solution}
$\cos x=2$ 的解集是空集（因为 $\cos x\leq1$），更不可能含有区间。
\end{solution}

\subsection{视角六：策略——先想哪一边好证}

\key{要点}
\begin{itemize}
  \item ∃ 通常好证：断言弱，而且有存在性定理白送。
  \item ∀ 难证：必须对所有情形给出一致论证，一个反例就推翻。
  \item ¬∃（不存在）通常最难：要排除所有可能。
  \item 所以先想「证哪一边便宜」，再决定方向。
\end{itemize}

\begin{example}
「方程 $x^3-3x+1=0$ 在 $(0,1)$ 内有解」与「在 $(0,1)$ 内无解」，哪个容易证？
\end{example}
\begin{solution}
有解容易：令 $f(x)=x^3-3x+1$，$f(0)=1>0$，$f(1)=-1<0$，零点定理直接给出。
无解才难：要说明解集为空，得排除所有可能。
\key{（先看结构，再决定证哪一边）}
\end{solution}

\begin{problem}
说明「$\forall x\in\left(0,\dfrac{\pi}{2}\right)$，$\sin x<x$」的证明思路，
并指出它比「$\exists x\in\left(0,\dfrac{\pi}{2}\right)$，$\sin x<x$」难在哪里。
\end{problem}
\begin{solution}
$\exists$ 只需找一个值（如 $x=\dfrac{\pi}{6}$，$\sin\dfrac{\pi}{6}=\dfrac12<\dfrac{\pi}{6}$）；
$\forall$ 要对整段区间成立，需设 $f(x)=x-\sin x$，由 $f'(x)=1-\cos x\geq0$ 且 $f(0)=0$
得 $f(x)>0$，是一致论证。
\end{solution}
\end{explore}

\begin{inclasstest}
\subsection{变式训练}
只改一个条件，重新做一遍。
\begin{enumerate}
  \item 写出「$\forall x\in\mathbb{R}$，$x^2+x+1>0$」的否定；
  \item 把「$\exists x\in[1,3]$，$x^2-4x+3\leq0$」写成集合语言并判断；
  \item 证明：存在 $x\in(0,1)$，使 $x^2+x-1=0$；
  \item 用反证法证明：对任意 $x\in\mathbb{R}$，$x^2-2x+3>0$；
  \item 证明：长度总和小于 $2$ 的有限个开区间，不能覆盖 $[0,2]$；
  \item 证明：不存在长度大于 $0$ 的区间 $I$，使 $\cos x=\cos y$ 对所有 $x,y\in I$ 成立。
\end{enumerate}
\end{inclasstest}

\begin{afterclass}
\subsection{知识来源对照表}
\begin{center}
\footnotesize
\setlength{\tabcolsep}{4pt}
\renewcommand{\arraystretch}{1.35}
\begin{tabular}{p{3.4cm}p{4.0cm}p{5.6cm}}
\toprule
结论 / 方法 & 最早从哪来 & 后面从哪个角度看 \\
\midrule
∃ 与 ∀ 的互化 & 视角一（否定） & 视角四（反证法） \\
解集与量词对应 & 视角二（结构） & 视角一（否定） \\
存在性定理 & 视角三（存在） & 视角六（策略） \\
反证法 & 视角四（反证） & 视角一（否定） \\
长度、距离、个数 & 视角五（度量） & 视角二（结构） \\
先想哪一边好证 & 视角六（策略） & 视角三（存在） \\
\bottomrule
\end{tabular}
\end{center}

\key{下一步}：量词是「集合语言」和「逻辑语言」的接口——
$\exists$ 对应解集非空，$\forall$ 对应解集等于定义域，$\neg$ 对应补集。
把这三件事记牢，恒成立、能成立、无解、存在性就都成了同一套语言。
\end{afterclass}
